For the direct equation, the first independent checker verifies 929 literal
conditions and rejects four deliberately altered artifacts: a changed
orientation term, a wrong first-factor shift, an escaping private quotient,
and a changed collected coefficient. A separately written reviewer checks
all 199 retained rows, all seven new rows, the nine actual raw power ports,
and all 216 parent restoration polynomials. None of those restorations
depends on the removed private guard coordinates or the changed private
pool quotient. The second reviewer re-expands the full quartic by unordered
monomial pairs, independently of the producer's ordered-pair algorithm.
The coefficient of $z_0^4$ is 85, a literal degree-attainment witness.

For each graph profile the checker first verifies the original normalized
source and graph certificate. It then checks the complete logical alias
map, separate flow coordinates, regenerated equality constraints and unit
path. It reconstructs twice the final polynomial as a sum of squares,
including the three-square identity for each quartic pair guard, and
compares every collected coefficient. It also parses every term of the
readable \path{polynomial.txt} back to the normative coefficient list.

All three graph profiles have been reconstructed from their source rows,
with byte-identical parent polynomials, parent certificates, final
polynomials, final certificates, and metric files: eighteen mathematical
JSON artifacts in total. The optional reconstruction uses the builders only
after the independent verification. Normal and optimization-mode Python
runs retain all necessary checks; the verifiers do not depend on
\texttt{assert} statements.
The integrated reconstruction also recreates all six direct-equation files
in a temporary directory and compares their exact bytes.

The graph evaluator also supplies a canonical witness lift. For an arbitrary
integer source tuple $z$, its construction satisfies the exact identity
\[
 2P_{\mathrm{unit}}(L(z))
   =\sum_rR_r(z)^2+\left(\sum_rR_r(z)\right)^2.
\]
On source zeros this becomes a constructive lift to graph zeros. The
identity follows from the block-tree construction and the monic coefficient
normalization; finite evaluations are supplementary diagnostics.

The integrated release check additionally compares the 306-variable source
copied into the graph component with the direct equation byte for byte,
parses every printed LaTeX residual back into its exact coefficient map,
and checks the equation identities and integrity inventory. Its receipt is
\path{verification/RELEASE_CHECK.json}; the optimization-mode companion is
\path{verification/RELEASE_CHECK_OPTIMIZED.json}. Neither check is a search
for a zero or a finite test of consistency.
The retained build receipts precede final inventory sealing; running the
released checker also verifies the completed \path{MANIFEST.sha256}.
